%========================================
% MatinBook — Stage 04: Typography Test
% Version: v1.1
% Purpose:
%   Validate the advanced typography system of MatinBook v1.1:
%     1. Microtype (protrusion, expansion)
%     2. Line breaking and hyphenation
%     3. Widow and orphan control
%     4. Special characters (quotes, dashes, brackets)
%     5. Latin typography and ligatures
%     6. Mixed Persian/Latin paragraphs
%     7. Mathematical spacing
%     8. Inline code typography
%     9. Draft mode watermark
%
% Compilation:
%   xelatex -shell-escape stage04-typography.tex
%   xelatex -shell-escape stage04-typography.tex
%   (2 passes required for cover + cross-references)
%
% Expected outcome:
%   A professionally typeset PDF demonstrating that all
%   typography settings are applied correctly, with proper
%   line breaking, spacing, and margin protrusion.
%========================================

%------------------------------------------------------------------------
% Search Path (for tests directory)
%------------------------------------------------------------------------
\makeatletter
\def\input@path{{../}}
\makeatother

%------------------------------------------------------------------------
% Document Class
%------------------------------------------------------------------------
\documentclass{matinbook}

%------------------------------------------------------------------------
% Book Metadata
%------------------------------------------------------------------------
\title{آزمون تایپوگرافی پیشرفته}
\author{تیم توسعه MatinBook}
\date{مهر ۱۴۰۵}

\booktitle{آزمون تایپوگرافی پیشرفته}

\renewcommand{\repository}{github.com/matinmahmoudi-cs/matinbook}

%========================================================================
% DOCUMENT
%========================================================================
\begin{document}

%------------------------------------------------------------------------
% FRONT COVER
%------------------------------------------------------------------------
\makecover
    {آزمون تایپوگرافی پیشرفته}
    {ارزیابی خط‌شکنی، فاصله‌گذاری و ریزحروف‌چینی}
    {تیم توسعه MatinBook}
    {مهر ۱۴۰۵}

%------------------------------------------------------------------------
% FRONT MATTER (symmetric margins)
%------------------------------------------------------------------------
\frontmatter
\frontmattergeometry

% Title page
\maketitle

% Copyright / Colophon
\thispagestyle{empty}
\vfill
\begin{center}
    {\large\textbf{شناسنامه آزمون}}

    \vspace{0.8cm}

    \begin{tabular}{rl}
        عنوان: & آزمون تایپوگرافی پیشرفته \\
        نسخه: & \matinbookversion \\
        تاریخ: & \matinbookdate \\
        موتور: & \lr{XeLaTeX} \\
        مجوز: & \lr{MIT} \\
    \end{tabular}

    \vspace{0.8cm}

    {\small این سند بخشی از مجموعه آزمون‌های یکپارچگی \lr{MatinBook} است.}
\end{center}
\clearpage

% Preface
\chapter*{پیشگفتار}
\addcontentsline{toc}{chapter}{پیشگفتار}

این سند، چهارمین آزمون از مجموعه‌ی آزمون‌های یکپارچگی
\lr{MatinBook} نسخه‌ی \lr{\matinbookversion} است. هدف آن،
اعتبارسنجی کامل سیستم تایپوگرافی کلاس در نه حوزه‌ی متمایز
است: ریزحروف‌چینی (\lr{microtype})، خط‌شکنی و هایفن‌گذاری،
کنترل خطوط تنها، نویسه‌های خاص، تایپوگرافی لاتین، متن
ترکیبی، فاصله‌گذاری ریاضی، کد درون‌خطی و حالت پیش‌نویس.

تایپوگرافی حرفه‌ای، تفاوت میان یک کتاب معمولی و یک کتاب
فاخر است. تنظیمات دقیق فاصله‌گذاری، بیرون‌زدگی نویسه‌ها
و مدیریت خطوط تنها، همگی به خوانایی و زیبایی متن کمک
می‌کنند. \lr{MatinBook} با استفاده از \lr{microtype} و
تنظیمات سفارشی، این اهداف را محقق می‌سازد.

هر بخش از این آزمون، با مثال‌های عملی و ارجاعات متقابل
همراه است تا خواننده بتواند درستی رفتار هر زیرسیستم را
در بافت واقعی یک کتاب ارزیابی کند.

\clearpage

% Table of Contents
\tableofcontents
\clearpage

% List of Figures
\listoffigures
\clearpage

% List of Tables
\listoftables
\clearpage

%------------------------------------------------------------------------
% MAIN MATTER (wide margin for notes)
%------------------------------------------------------------------------
\mainmatter
\mainmattergeometry

%========================================================================
% CHAPTER 1 — PERSIAN TYPOGRAPHY
%========================================================================
\chapter{تایپوگرافی فارسی}
\label{chap:typography-persian}

\section{فلسفه‌ی تایپوگرافی در MatinBook}
\label{sec:typography-philosophy}

تایپوگرافی خوب، هنر و فن چیدمان حروف و کلمات برای ایجاد
متنی خوانا، شفاف و جذاب است. \lr{MatinBook} با استفاده
از بسته‌ی \lr{microtype} و تنظیمات دقیق خط‌شکنی
(\lr{line breaking})، تجربه‌ی خواندن متون فارسی را بهبود
می‌بخشد.

تنظیمات \lr{microtype} در \lr{XeTeX} شامل سه بخش اصلی است:

\begin{itemize}
    \item \textbf{بیرون‌زدگی (\lr{Protrusion}):} فعال — کاراکترهای خاص
    مانند نقطه، کاما و گیومه کمی از حاشیه بیرون می‌زنند
    تا لبه‌ی متن صاف‌تر به نظر برسد.
    \item \textbf{انبساط (\lr{Expansion}):} غیرفعال — \lr{XeTeX}
    از این قابلیت پشتیبانی نمی‌کند.
    \item \textbf{کرنینگ (\lr{Kerning}):} توسط خود قلم مدیریت
    می‌شود و به \lr{microtype} واگذار نمی‌شود.
\end{itemize}

\mnote{بیرون‌زدگی نویسه‌ها، لبه‌ی متن را بصری صاف‌تر می‌کند، اگرچه در واقعیت نویسه‌ها کمی از حاشیه بیرون می‌زنند.}

\section{متن طولانی فارسی}
\label{sec:typography-longtext}

این یک پاراگراف بسیار طولانی است که به‌طور اختصاصی برای
آزمایش نحوه‌ی شکسته شدن خطوط در متن فارسی نوشته شده است.
هدف ما این است که ببینیم آیا کلمات به‌درستی در انتهای
خطوط می‌شکنند و توزیع فضای خالی بین کلمات به‌صورت یکنواخت
انجام می‌شود یا خیر.

در حروف‌چینی فارسی، برخلاف انگلیسی، کلمات معمولاً شکسته
نمی‌شوند مگر در موارد خاص. بنابراین تنظیمات خط‌شکنی باید
به گونه‌ای باشد که از ایجاد فضاهای خیلی بزرگ یا خیلی کوچک
بین کلمات جلوگیری کند. \lr{MatinBook} با تنظیم
\lr{\textbackslash emergencystretch=1em} و
\lr{\textbackslash tolerance=3000} این مشکل را مدیریت
می‌کند.

\section{جلوگیری از خطوط تنها}
\label{sec:typography-widow}

در حروف‌چینی حرفه‌ای فارسی، دو مشکل رایج وجود دارد:

\begin{enumerate}
    \item \textbf{خط بی‌آویز (\lr{Widow}):} آخرین خط یک پاراگراف
    که به‌تنهایی در ابتدای صفحه‌ی بعد قرار می‌گیرد.
    \item \textbf{خط یتیم (\lr{Orphan}):} اولین خط یک پاراگراف
    که به‌تنهایی در انتهای صفحه‌ی قبل باقی می‌ماند.
\end{enumerate}

\lr{MatinBook} با تنظیم \lr{\textbackslash widowpenalty=10000}
و \lr{\textbackslash clubpenalty=10000} از هر دو مشکل
جلوگیری می‌کند. این مقادیر بالا به \lr{LaTeX} می‌گویند
که هرگز اجازه‌ی ایجاد این خطوط را ندهد.

\section{نویسه‌های خاص در تایپوگرافی}
\label{sec:typography-special}

نویسه‌های خاص فارسی و نشانه‌گذاری:

\begin{itemize}
    \item \textbf{گیومه:} «متن داخل گیومه» باید چسبیده به متن باشد.
    \item \textbf{پرانتز:} (متن داخل پرانتز) بدون فاصله‌ی اضافی.
    \item \textbf{نقطه:} پایان جمله. فاصله بعد از نقطه باید
    استاندارد باشد.
    \item \textbf{کاما:} برای جداسازی، عبارات داخل جمله.
    \item \textbf{خط تیره:} برای عبارت‌های معترضه — مانند
    این — استفاده می‌شود.
\end{itemize}

%========================================================================
% CHAPTER 2 — LATIN TYPOGRAPHY
%========================================================================
\chapter{تایپوگرافی لاتین}
\label{chap:typography-latin}

\section{تنظیمات microtype}
\label{sec:typography-microtype-latin}

\begin{latin}
This chapter demonstrates advanced Latin typography features
provided by the \texttt{microtype} package in MatinBook.

\textbf{Protrusion Settings.}
The following characters are configured to protrude slightly
into the margins, creating a visually straight edge:

\begin{itemize}
    \item Period (.) --- protrudes by 700/1000
    \item Comma (,) --- protrudes by 700/1000
    \item Colon (:) --- protrudes by 500/1000
    \item Semicolon (;) --- protrudes by 500/1000
    \item Quotation marks --- protrude by 1000/1000
\end{itemize}
\end{latin}

\section{متن طولانی لاتین}
\label{sec:typography-longtext-latin}

\begin{latin}
This is a deliberately long paragraph designed to test the line
breaking algorithm in English text. The quick brown fox jumps
over the lazy dog multiple times to ensure that the text fills
enough lines to demonstrate proper justification behavior. We
need to verify that hyphenation works correctly when necessary,
and that the spacing between words remains consistent
throughout the paragraph without creating large gaps or rivers
of white space that would distract the reader from the content.
\end{latin}

\section{هایفن‌گذاری}
\label{sec:typography-hyphenation}

\begin{latin}
The following words should hyphenate correctly when near line
breaks: implementation, configuration, initialization,
documentation, communication, functionality, comprehensive,
sophisticated, characteristic, representation.
\end{latin}

\section{لیگچرها}
\label{sec:typography-ligatures}

\begin{latin}
Ligatures are special characters that combine two or more
letters into a single glyph. Common ligatures in English
include:

\begin{itemize}
    \item fi ligature: define, finance, profile
    \item fl ligature: flower, reflect, inflate
    \item ff ligature: offer, different, difficult
    \item ffi ligature: efficient, coefficient
    \item ffl ligature: shuffle, waffle, baffle
\end{itemize}

Times New Roman supports all standard ligatures automatically.
\end{latin}

%========================================================================
% CHAPTER 3 — MIXED TEXT TYPOGRAPHY
%========================================================================
\chapter{تایپوگرافی ترکیبی}
\label{chap:typography-mixed}

\section{فارسی و انگلیسی در یک پاراگراف}
\label{sec:typography-mixed-paragraph}

نوشتن متون علمی-فنی به زبان فارسی همواره با چالش ترکیب
واژگان انگلیسی همراه است. برای مثال، وقتی درباره
\lr{machine learning} صحبت می‌کنیم، باید مفاهیمی مانند
\lr{neural network}، \lr{deep learning}، \lr{gradient descent}
و \lr{backpropagation} را در متن فارسی بگنجانیم.
همچنین اسامی خاص مانند \lr{TensorFlow}، \lr{PyTorch}،
\lr{Scikit-learn} و \lr{Jupyter Notebook} نیز به وفور
در متون برنامه‌نویسی دیده می‌شوند.

\lr{MatinBook} با استفاده از \lr{xepersian} و دستور
\lr{\textbackslash lr} برای کلمات لاتین، این ترکیب را
به‌صورت حرفه‌ای مدیریت می‌کند. فاصله قبل و بعد از کلمات
لاتین به‌صورت خودکار تنظیم می‌شود.

\section{فرمول‌های ریاضی در متن}
\label{sec:typography-math-inline}

فرمول‌های ریاضی نیز بخش جدایی‌ناپذیر متون علمی هستند.
معادلات درون‌خطی مانند $f(x) = x^{2} + 3x - 1$ یا
$\sum_{i=1}^{n} i = \frac{n(n+1)}{2}$ باید بدون برهم
زدن فاصله‌ی خطوط در متن جای گیرند.

فرمول‌های نمایشی نیز باید با فاصله‌ی مناسب از متن قبل و
بعد قرار گیرند:

\[
    \mathcal{F}(\omega) = \int_{-\infty}^{\infty} f(t) \, e^{-i\omega t} \, dt
\]

\section{کدهای درون‌خطی}
\label{sec:typography-inline-code}

گاهی لازم است به توابع یا متغیرها در متن اشاره کنیم.
برای مثال، تابع \inlcode{calculate_statistics()} یک لیست
از اعداد را دریافت می‌کند، یا متغیر \inlcode{total_count}
تعداد کل را نگه می‌دارد. این ارجاعات باید با فونت
\lr{monospace} و بدون برهم زدن نظم پاراگراف نمایش داده شوند.

%========================================================================
% CHAPTER 4 — MATHEMATICAL SPACING
%========================================================================
\chapter{فاصله‌گذاری ریاضی}
\label{chap:typography-math}

\section{تنظیمات فاصله در مد ریاضی}
\label{sec:typography-math-spacing}

\begin{latin}
\LaTeX\ provides fine-grained control over spacing in
mathematical mode. MatinBook configures these parameters for
optimal appearance:

\begin{itemize}
    \item \texttt{\textbackslash thinmuskip} = 3mu --- thin space
    (e.g., around punctuation in math)
    \item \texttt{\textbackslash medmuskip} = 4mu plus 2mu minus 4mu
    --- medium space (e.g., around binary operators)
    \item \texttt{\textbackslash thickmuskip} = 5mu plus 5mu
    --- thick space (e.g., around relational operators)
\end{itemize}
\end{latin}

\section{مقایسه فاصله‌ها}
\label{sec:typography-math-comparison}

\begin{align}
    a + b &= c + d \label{eq:spacing1} \\
    a+b &= c+d \label{eq:spacing2} \\
    \int f(x) \, dx &\neq \int f(x) dx \label{eq:spacing3}
\end{align}

توجه کنید که در معادله‌ی \cref{eq:spacing1}، فاصله دور
عملگر $+$ بیشتر از معادله‌ی \cref{eq:spacing2} است.
همچنین در معادله‌ی \cref{eq:spacing3}، استفاده از
\lr{\textbackslash ,} قبل از $dx$ فاصله‌ی استاندارد
را ایجاد می‌کند.

\section{شکستن معادلات طولانی}
\label{sec:typography-math-breaks}

\lr{MatinBook} با فعال‌سازی
\lr{\textbackslash allowdisplaybreaks} اجازه می‌دهد
معادلات طولانی در محیط‌های \lr{align} و \lr{gather}
بین صفحات بشکنند. این ویژگی برای کتاب‌های ریاضی که
شامل اثبات‌های طولانی هستند، ضروری است.

%========================================================================
% CHAPTER 5 — DRAFT MODE
%========================================================================
\chapter{حالت پیش‌نویس}
\label{chap:typography-draft}

\section{تنظیمات Draft}
\label{sec:typography-draft-settings}

در حالت \lr{draft}، \lr{MatinBook} رفتارهای زیر را تغییر
می‌دهد:

\begin{itemize}
    \item \textbf{واترمارک:} عبارت \lr{DRAFT} با شفافیت
    \pnum{۹۰٪} در پس‌زمینه‌ی صفحات نمایش داده می‌شود.
    \item \textbf{جعبه‌های سرریز:} خطوطی که از حاشیه خارج
    می‌شوند با خط مشکی \pnum{۵pt} مشخص می‌شوند.
    \item \textbf{تصاویر:} به‌جای نمایش تصاویر، فقط
    کادر آن‌ها رسم می‌شود (سرعت کامپایل بالاتر).
\end{itemize}

برای فعال‌سازی حالت \lr{draft}، از دستور زیر استفاده کنید:

\begin{latin}
\begin{minted}{latex}
\documentclass[draft]{matinbook}
\end{minted}
\end{latin}

\begin{remark}
    حالت \lr{draft} برای مراحل اولیه‌ی نوشتن کتاب بسیار
    مفید است، زیرا سرعت کامپایل را افزایش می‌دهد و مشکلات
    صفحه‌آرایی را برجسته می‌کند. قبل از چاپ نهایی، حتماً
    به حالت \lr{final} برگردید.
\end{remark}

%========================================================================
% CHAPTER 6 — SUMMARY
%========================================================================
\chapter{خلاصه‌ی نتایج}
\label{chap:summary}

\section{وضعیت تایپوگرافی}
\label{sec:summary-typography}

جدول \cref{tab:typography-status} وضعیت تنظیمات تایپوگرافی
\lr{MatinBook} را نشان می‌دهد.

\begin{matintable}{وضعیت تنظیمات تایپوگرافی}{tab:typography-status}
\begin{tblr}{
    width = \textwidth,
    colspec = {l X[c] X[c]},
    hlines, vlines,
    colsep = 8pt,
    rowsep = 4pt,
    row{1} = {font=\bfseries},
}
تنظیم & وضعیت & توضیحات \\
\lr{Microtype Protrusion} & \textcolor{matingreen}{فعال} & بیرون‌زدگی نویسه‌های خاص \\
\lr{Microtype Expansion} & \textcolor{matinorange}{غیرفعال} & \lr{XeTeX} پشتیبانی نمی‌کند \\
\lr{Widow Penalty} & \textcolor{matingreen}{\pnum{۱۰۰۰۰}} & جلوگیری از خط بی‌آویز \\
\lr{Club Penalty} & \textcolor{matingreen}{\pnum{۱۰۰۰۰}} & جلوگیری از خط یتیم \\
\lr{Emergency Stretch} & \textcolor{matingreen}{\pnum{۱em}} & فضای اضطراری خط‌شکنی \\
\lr{Overfull Rule} & \textcolor{matingreen}{\pnum{۰pt}} & مخفی در حالت نهایی \\
\lr{Math Spacing} & \textcolor{matingreen}{بهینه} & thin/med/thick muskip \\
\lr{Display Breaks} & \textcolor{matingreen}{مجاز} & شکست معادلات طولانی \\
\end{tblr}
\end{matintable}

\section{معیارهای ارزیابی}
\label{sec:summary-criteria}

در این آزمون، سیستم تایپوگرافی \lr{MatinBook} بر اساس
چهار معیار ارزیابی شد:

\begin{enumerate}
    \item \textbf{خوانایی}: در متن‌های طولانی فارسی و لاتین.
    \item \textbf{ترازبندی}: کیفیت لبه‌های متن در دو طرف.
    \item \textbf{فاصله‌گذاری}: در متن، ریاضیات و کد درون‌خطی.
    \item \textbf{پایداری}: رفتار یکنواخت در حالت‌های مختلف.
\end{enumerate}

\section{نتیجه‌گیری}
\label{sec:summary-conclusion}

\textbf{تمام زیرسیستم‌های تایپوگرافی \lr{MatinBook} نسخه‌ی
\lr{\matinbookversion} با موفقیت بارگذاری و آزمایش شدند.
تنظیمات خط‌شکنی، فاصله‌گذاری و \lr{microtype} به‌درستی
کار می‌کنند و متن‌های فارسی و لاتین، خوانایی و زیبایی
مطلوبی دارند.}

%------------------------------------------------------------------------
% BACK MATTER (symmetric margins)
%------------------------------------------------------------------------
\backmatter
\frontmattergeometry

% Bibliography (empty placeholder)
\printbibliography[title={منابع و مراجع}]

% Index
\printindex

%------------------------------------------------------------------------
% BACK COVER
%------------------------------------------------------------------------
\authorbio{%
    تیم توسعه‌ی \lr{MatinBook}، گروهی از پژوهشگران و برنامه‌نویسان
    فارسی‌زبان است که با هدف ارائه‌ی یک راه‌حل حرفه‌ای برای
    حروف‌چینی کتاب‌های فنی فارسی فعالیت می‌کند. این پروژه حاصل
    سال‌ها تجربه در حوزه‌ی نشر دیجیتال و برنامه‌نویسی است.
}

\makebackcover{%
    این سند، چهارمین آزمون از مجموعه‌ی آزمون‌های یکپارچگی
    \lr{MatinBook} نسخه‌ی \lr{\matinbookversion} است.

    \vspace{0.5cm}

    \textbf{زیرسیستم‌های آزموده‌شده:}
    \begin{itemize}
        \item ریزحروف‌چینی \lr{microtype}
        \item خط‌شکنی و هایفن‌گذاری
        \item کنترل خطوط تنها
        \item نویسه‌های خاص و نگارشی
        \item تایپوگرافی لاتین و لیگچرها
        \item متن ترکیبی فارسی و لاتین
        \item فاصله‌گذاری ریاضی
        \item حالت پیش‌نویس
    \end{itemize}
}

\end{document}